\documentclass[11pt,a4paper]{article}

\usepackage[utf8]{inputenc}
\usepackage[T1]{fontenc}
\usepackage[english]{babel}
\usepackage{lmodern}
\usepackage{microtype}
\usepackage{geometry}
\geometry{margin=1in}
\usepackage{booktabs}
\usepackage{array}
\usepackage{enumitem}
\usepackage{hyperref}
\hypersetup{
    colorlinks=true,
    linkcolor=black,
    urlcolor=blue,
    citecolor=black,
    pdftitle={The Honesty Matrix: A Stipulative Two-Dimensional Typology of Truthful and Untruthful Utterances},
    pdfauthor={Norbert Nopper},
    pdfsubject={Philosophy of language; epistemology; typology of assertion}
}

\title{The Honesty Matrix:\\A Stipulative Two-Dimensional Typology\\of Truthful and Untruthful Utterances}
\author{Norbert Nopper}
\date{Version 1.0 --- 14 May 2026}

\begin{document}

\maketitle

\begin{abstract}
\noindent
We propose the \emph{Honesty Matrix}, a compact $2\times 2$ typology that
classifies declarative utterances along two independent axes: (i) the
\emph{truth-value} of the utterance with respect to the actual state of
affairs, and (ii) the speaker's \emph{awareness} of the full epistemic and
consequential status of that utterance at the moment of speaking. The
model is presented as a \emph{stipulative} framework, not as a
redefinition of the philosophical notion of lying. It is intended to
organise familiar phenomena --- sincere truth-telling, deliberate
deception, the Freudian slip, the white lie, the honest mistake, and
Frankfurtian bullshit --- inside a single, explicit structure. We give
precise working definitions, place the framework relative to existing
accounts in philosophy of lying and psychoanalysis, work each cell
through with examples, sketch a minimal operationalisation, and state
the model's limitations.
\end{abstract}

\noindent\textbf{Keywords:} honesty, lying, deception, sincerity,
accuracy, consciousness, self-deception, Freudian slip, bullshit,
typology.

\section{Introduction}

Standard philosophical analyses of lying require three elements: the
speaker makes a statement, believes it to be false, and intends to
deceive the hearer~\cite{augustine,bok1978,mahon2016sep,carson2010}. This
belief-centred account captures the paradigmatic lie but leaves several
familiar phenomena underdescribed: utterances whose truth-value the
speaker has not consciously considered (e.g.\ slips of the
tongue~\cite{freud1901}); utterances whose content is technically true
but produced without deliberate intent; ``socially benign'' falsehoods
whose downstream consequences are unknown~\cite{bok1978}; assertions
produced with indifference to truth~\cite{frankfurt2005}; and deceptive
communication achieved without uttering any literal
falsehood~\cite{saul2012,grice1975}.

This paper does not contest the standard definition of lying. Instead,
it introduces a \emph{descriptive typology of utterances} --- the
Honesty Matrix --- that adds an explicit \emph{awareness} dimension to
the truth-value axis already implicit in the literature. The matrix is
stipulative: its cells are defined by the axes, not by appeal to
ordinary usage. Its purpose is taxonomic clarity rather than conceptual
revision.

The closest precursor is Williams's distinction between the two
``virtues of truth'': \emph{Sincerity} (saying what one believes) and
\emph{Accuracy} (taking due care that one's beliefs are
true)~\cite{williams2002}. The Honesty Matrix can be read as a
discretised cousin of that view: the content axis tracks accuracy of
the utterance, while the awareness axis tracks something closely
related to --- though broader than --- sincerity. Related contemporary
treatments include Carson's~\cite{carson2010} warranting account of
lying, Sorensen's~\cite{sorensen2007} discussion of bald-faced lies,
and Stokke's~\cite{stokke2018} insincerity-based analysis. The present
contribution is the explicit $2\times 2$ tabulation, the
acknowledgement of an \emph{unconscious} mode of utterance on the
speaker's side, and the mapping of familiar phenomena onto the
resulting four cells.

\section{Definitions}

We adopt the following working definitions. They are stipulative within
this paper; alternative definitions exist and are discussed in
Section~\ref{sec:limits}.

\begin{description}
  \item[\textbf{Utterance.}] A declarative statement made by a speaker
        with assertoric force.
  \item[\textbf{Truth-value (content axis).}] An utterance is
        \emph{true} if it corresponds to the actual state of affairs,
        and \emph{false} otherwise. We use ``lie'' in this paper as a
        short label for ``false utterance'' --- a deliberate
        terminological simplification that diverges from the
        philosophical mainstream (cf.\ Section~\ref{sec:limits}).
        Correspondence is assumed only as a working assumption.
  \item[\textbf{Awareness (epistemic axis).}] A speaker is
        \emph{conscious} of an utterance when, at the moment of
        speaking, they have deliberate access to (a) its intended
        content, (b) its truth-value relative to their own beliefs,
        and (c) its foreseeable consequences. If any of these is
        absent, the speaker is, to that extent, \emph{unconscious} of
        the utterance.
\end{description}

\noindent
Two notes on the awareness axis. First, ``unconscious'' is used in a
broad psychological sense (lack of deliberate access), not strictly the
Freudian technical sense; it is closer to Block's notion of failure of
\emph{access consciousness}~\cite{block1995}. Second, awareness is
treated as binary here for classificatory simplicity; a graded version
is a natural extension and is briefly discussed in
Section~\ref{sec:limits}.

\section{The Matrix}

Crossing the two axes yields the four cells in Table~\ref{tab:matrix}.
Each cell name is a label internal to this paper; correspondences to
ordinary language are discussed in Section~\ref{sec:phenomena} and
illustrated by examples in Section~\ref{sec:examples}.

\begin{table}[h]
\centering
\caption{The Honesty Matrix.}
\label{tab:matrix}
\begin{tabular}{lll}
\toprule
            & \textbf{Conscious}            & \textbf{Unconscious} \\
\midrule
\textbf{True}  & Sincere truth-telling      & Unintended truth \\
\textbf{False} & Deliberate lie             & Unwitting falsehood \\
\bottomrule
\end{tabular}
\end{table}

\begin{description}
  \item[\textbf{Sincere truth-telling.}] The utterance is true, and the
        speaker has deliberate access to its content, truth-value, and
        likely consequences. This is the prototype of honesty.
  \item[\textbf{Deliberate lie.}] The utterance is false, and the
        speaker is aware of its falsity and of its intended deceptive
        function. This matches the standard philosophical definition of
        lying~\cite{mahon2016sep,carson2010}.
  \item[\textbf{Unintended truth.}] The utterance is true, but the
        speaker did not deliberately intend to assert it as such ---
        e.g.\ a slip of the tongue that happens to state a fact, or an
        offhand remark that turns out to be veridical.
  \item[\textbf{Unwitting falsehood.}] The utterance is false, but the
        speaker lacks full deliberate access to its falsity, its
        consequences, or both. The clearest occupants are sincere
        errors and certain forms of self-deception~\cite{mele2001}.
\end{description}

\section{Mapping Familiar Phenomena}
\label{sec:phenomena}

\subsection{Freudian slip $\to$ Unintended truth}
A parapraxis, in psychoanalytic theory, is a verbal error read as the
surfacing of unconscious material~\cite{freud1901}. Independently of
whether one accepts the psychoanalytic mechanism, the linguistic
structure of such slips fits the \emph{unintended truth} cell whenever
the slip happens to state something the speaker would not have
deliberately asserted but which is, in fact, the case. We make no claim
about the empirical frequency or psychological cause of such slips.

\subsection{Honest mistake and self-deception $\to$ Unwitting falsehood}
A speaker who sincerely asserts a falsehood they believe to be true is
classified here as producing an \emph{unwitting falsehood}. The same
cell hosts cases of self-deception, in which a speaker sincerely
asserts what, on some deeper level, they have reason to doubt or
disbelieve~\cite{mele2001}.

\subsection{The white lie: a boundary case}
The \emph{white lie} is conventionally a falsehood the speaker knows to
be false but believes to be socially harmless~\cite{bok1978}. Under the
definitions adopted here it sits \emph{between} the deliberate lie (the
speaker is aware of the falsity) and the unwitting falsehood (the
speaker is typically unaware of the lie's full long-term consequences).
It is therefore best treated as a boundary case rather than a clean
occupant of either cell. The matrix does not aim to dissolve this
tension; it makes it explicit.

\subsection{Bullshit: a phenomenon the matrix does not capture}
Frankfurt's~\cite{frankfurt2005} bullshitter is characterised by
\emph{indifference} to truth-value, not by unawareness of it. Bullshit
is therefore not a clean occupant of any cell of the present matrix: it
represents a third epistemic stance --- neither believing-true nor
believing-false --- and would require a richer awareness axis (or an
explicit ``indifference'' attitude) to be modelled directly. We note
this as a known incompleteness of the model.

\subsection{Deception without lying}
A speaker can deceive while uttering only true sentences --- through
conversational implicature, selective emphasis, misleading omission, or
strategic silence~\cite{grice1975,saul2012}. Such cases lie outside the
scope of the matrix: the literal-content axis cannot capture them. The
matrix is best read as covering \emph{assertoric} dishonesty, not the
wider category of \emph{deceptive communication}.

\subsection{Irony, fiction, and performatives}
Ironic, fictional, and performative utterances have literal contents
whose truth-value diverges from what is asserted. We follow standard
practice~\cite{grice1975} in treating these as cases in which
assertoric force is suspended or shifted; they fall outside the matrix
because no straightforward truth-value evaluation of the literal
content applies.

\section{Worked Examples}
\label{sec:examples}

The following four examples instantiate one occupant of each cell. They
are illustrative, not empirical.

\begin{description}[style=nextline]
  \item[\textbf{Sincere truth-telling (True, Conscious).}]
    A witness on the stand, having clearly recalled an event and
    intending to convey what they remember, states, ``The car was
    red.'' The car was, in fact, red.
  \item[\textbf{Deliberate lie (False, Conscious).}]
    A defendant, knowing he was at the scene, states, ``I was at home
    that evening,'' intending the court to believe him.
  \item[\textbf{Unintended truth (True, Unconscious).}]
    A spouse means to say ``I'll be at the office'' but accidentally
    says ``I'll be at the bar.'' They are, in fact, going to the bar,
    though they did not consciously intend to disclose this.
  \item[\textbf{Unwitting falsehood (False, Unconscious).}]
    A speaker confidently asserts ``The capital of Australia is
    Sydney,'' sincerely believing it, unaware that the capital is
    Canberra.
\end{description}

\section{Operationalisation}

The matrix is conceptual, but a minimum procedure for assigning a given
utterance $u$ produced by speaker $S$ to a cell can be stated:

\begin{enumerate}
  \item Determine the truth-value of $u$ relative to the actual state
        of affairs. If unavailable, cell assignment along the content
        axis is suspended and the matrix is silent on $u$ until
        further evidence is obtained.
  \item Determine, from $S$'s subsequent reports, behaviour, and
        background beliefs, whether $S$ had deliberate access to (a)
        the intended content of $u$, (b) its truth-value relative to
        $S$'s own beliefs, and (c) its foreseeable consequences. Full
        deliberate access on all three places $u$ in the conscious
        row; failure on any one places $u$ in the unconscious row.
  \item If $u$ is non-assertoric (ironic, fictional, performative) or
        achieves its communicative effect via implicature rather than
        literal content, mark $u$ as out of scope.
\end{enumerate}

\noindent
The procedure is fallible: step~2 depends on testimony or inference
about a speaker's mental state, which is rarely decisive, and any
graded notion of awareness will yield borderline cases. The procedure
nevertheless makes the matrix's cell assignments in principle
checkable rather than purely declarative.

\section{Discussion}

Three observations follow from the matrix.

\paragraph{Honesty is not one-dimensional.} Two utterances with
identical truth-value can occupy different cells depending on the
speaker's awareness; two utterances at the same awareness level can
differ in truth-value. Any single-axis evaluation of ``honesty''
therefore omits information.

\paragraph{Truth can be incidental.} The \emph{unintended truth} cell
shows that veridical speech is not, by itself, evidence of deliberate
honesty.

\paragraph{Sincere assertion does not guarantee truth.} The
\emph{unwitting falsehood} cell shows the converse: deliberate honesty
is not evidence of veridical speech. Both observations are consistent
with the standard epistemological distinction between knowledge and
justified belief, and with knowledge-norm accounts of
assertion~\cite{williamson2000}.

\section{Limitations and Scope}
\label{sec:limits}

The matrix is deliberately narrow. The following limitations apply.

\begin{enumerate}
  \item \textbf{Stipulative definitions.} The use of ``lie'' for any
        false utterance (rather than a belief-and-intent-laden notion)
        is a terminological choice for this paper. Readers preferring
        the standard definition can substitute ``false utterance'' for
        ``lie'' throughout without loss.
  \item \textbf{Binary axes.} Truth-value and awareness are both
        treated as binary. Real cases are graded (partial truths,
        partial awareness); the binary form is a simplification.
        Graded axes would yield a continuous version of the
        matrix~\cite{block1995}.
  \item \textbf{Correspondence assumption.} The content axis assumes a
        correspondence-style notion of truth. Coherence- or
        pragmatist-style accounts would reshape the content axis but
        leave the awareness axis intact.
  \item \textbf{Partial axis coupling.} The axes are presented as
        independent, but ``deliberate lie'' presupposes awareness of
        falsity by definition. The independence claim is therefore
        descriptive, not strictly logical.
  \item \textbf{Assertoric scope.} The matrix covers declarative
        assertions only. Lying by omission, misleading implicature,
        irony, fiction, and performative utterances fall outside its
        scope~\cite{saul2012,grice1975}.
  \item \textbf{Indifference not modelled.} Frankfurtian
        bullshit~\cite{frankfurt2005} --- assertion without regard to
        truth --- is not captured by a binary awareness axis.
  \item \textbf{No empirical claims.} The mapping of Freudian slips,
        white lies, and self-deception to particular cells is
        conceptual, not empirical. No claim is made about how often
        such utterances occur or about their underlying psychological
        mechanisms.
  \item \textbf{Moral neutrality.} The matrix classifies utterances;
        it does not assign blame or praise. Moral evaluation typically
        depends on awareness, but is left to ethical theory.
\end{enumerate}

\section{Conclusion}

The Honesty Matrix is a small, explicit, stipulative typology. It
crosses the truth-value of an utterance with the speaker's awareness of
that utterance's status and identifies four cells: \emph{sincere
truth-telling}, \emph{unintended truth}, \emph{deliberate lie}, and
\emph{unwitting falsehood}. Familiar phenomena --- the Freudian slip,
the honest mistake, the white lie, and Frankfurtian bullshit --- can be
located, partially located, or explicitly recognised as out of scope
within this structure. The framework is offered as an organising tool,
compatible with rather than competing against established accounts of
lying, deception, self-deception, and parapraxis, and its limitations
are stated openly.

\begin{thebibliography}{99}

\bibitem{augustine}
Augustine of Hippo (c.~395).
\emph{De Mendacio} [\emph{On Lying}].
Reprinted in: \emph{Treatises on Various Subjects}, R. J. Deferrari
(ed.), Catholic University of America Press, 1952.

\bibitem{block1995}
Block, N. (1995).
On a confusion about a function of consciousness.
\emph{Behavioral and Brain Sciences}, 18(2), 227--247.

\bibitem{bok1978}
Bok, S. (1978).
\emph{Lying: Moral Choice in Public and Private Life.}
Pantheon Books, New York.

\bibitem{carson2010}
Carson, T. L. (2010).
\emph{Lying and Deception: Theory and Practice.}
Oxford University Press.

\bibitem{frankfurt2005}
Frankfurt, H. G. (2005).
\emph{On Bullshit.}
Princeton University Press.

\bibitem{freud1901}
Freud, S. (1901).
\emph{Zur Psychopathologie des Alltagslebens}
[\emph{The Psychopathology of Everyday Life}].
Karger, Berlin.

\bibitem{grice1975}
Grice, H. P. (1975).
Logic and conversation.
In: P. Cole \& J. L. Morgan (eds.),
\emph{Syntax and Semantics, Vol.~3: Speech Acts},
Academic Press, pp.~41--58.

\bibitem{mahon2016sep}
Mahon, J. E. (2016).
\emph{The Definition of Lying and Deception.}
The Stanford Encyclopedia of Philosophy (Winter 2016 ed.),
E. N. Zalta (ed.).
\url{https://plato.stanford.edu/archives/win2016/entries/lying-definition/}

\bibitem{mele2001}
Mele, A. R. (2001).
\emph{Self-Deception Unmasked.}
Princeton University Press.

\bibitem{saul2012}
Saul, J. M. (2012).
\emph{Lying, Misleading, and What is Said: An Exploration in Philosophy
of Language and in Ethics.}
Oxford University Press.

\bibitem{sorensen2007}
Sorensen, R. (2007).
Bald-faced lies! Lying without the intent to deceive.
\emph{Pacific Philosophical Quarterly}, 88(2), 251--264.

\bibitem{stokke2018}
Stokke, A. (2018).
\emph{Lying and Insincerity.}
Oxford University Press.

\bibitem{williams2002}
Williams, B. (2002).
\emph{Truth and Truthfulness: An Essay in Genealogy.}
Princeton University Press.

\bibitem{williamson2000}
Williamson, T. (2000).
\emph{Knowledge and Its Limits.}
Oxford University Press.

\end{thebibliography}

\end{document}
